\textcolor{blue}{2. [15 min/30 puntos] Considera que en una empresa necesitan computar un problema, sin
embargo solo se puede paralelizar el 80$\%$ de su cómputo. Al dueño de la empresa no le importa
invertir en tecnología. ¿De cuántos cores necesitaría comprar su computadora para aumentar la
velocidad de cómputo del problema al triple?}\\

A comparación del problema anterior aquí sólo se puede paralelizar el 80$\%$ por lo que de nuevo tendremos que despejar la Ley de Amdahl. Como datos tenemos 

$$p = 0.80$$
$$S=3$$

por lo que tenemos que encontrar $n$

$$S(n)=\frac{1}{(1-p)+\frac{p}{n}} $$

despejamos 

$$\frac{1}{S(n)}= (1-p)+\frac{p}{n}$$
$$\frac{1}{S(n)}-(1-p)=\frac{p}{n}$$
$$n\Big(\frac{1}{S(n)}-(1-p)\Big)=p$$
$$\frac{n(\frac{1}{S}-(1-p))}{\frac{1}{S}-(1-p)}= \frac{p}{\frac{1}{S}-(1-p)}$$
$$n = \frac{p}{\frac{1}{S}-(1-p)}$$

sustituimos

$$ n = \frac{0.80}{\frac{1}{3}-(1-0.80)}= 6$$

En conclusión la empresa necesitaría una computadora con 6 cores para obtener una aceleración al triple.\\